% =============================================================================
% 3. fejezet — Kutatási rész: a bőrelváltozás-osztályozó modell
% -----------------------------------------------------------------------------
% MÓDSZERTAN kész. Az EREDMÉNYEK szekció táblázatai/ábrái placeholderek,
% amelyeket a végleges v4 ViT és ResNet-18 futtatások kimenetéből kell kitölteni
% (metrics_final.json, training_curves.png, roc_curve_test.png).
% =============================================================================

Ez a fejezet a bőrelváltozás-osztályozó modell fejlesztését és kiértékelését
mutatja be. Először az ISIC 2024 adathalmaz előkészítését ismertetjük, majd a
Vision Transformer modellt és annak tanítási folyamatát, ezt követően a két
viszonyítási alapot (a ResNet-18 konvolúciós hálót és a Gemini nagy nyelvi
modellt), végül a kísérleti elrendezést és az elért eredményeket.

% =============================================================================
\section{Az ISIC 2024 adathalmaz és előkészítése}
% =============================================================================

A kísérletekhez az ISIC 2024 verseny (\textit{Skin Cancer Detection with 3D-TBP})
nyilvános tanító adathalmazát használtuk \cite{isic2024}. Az adathalmaz a teljes
testfelszín háromdimenziós fényképezéséből kivágott, mintegy 401~000
bőrelváltozás-felvételt tartalmaz, amelyek mindegyikéhez bináris címke tartozik
(jó\-indulatú vagy rosszindulatú). Az adathalmaz meghatározó tulajdonsága a
rendkívül erős osztály-egyensúlytalanság: a rosszindulatú esetek száma mindössze
393, ami megközelítőleg $1{:}1000$ pozitív--negatív arányt jelent.

A v3 előzetes kísérletektől eltérően -- amelyek a többségi (jóindulatú) osztály
50\%-os aldmintázásával dolgoztak -- a végleges modellt a \textbf{teljes}
adathalmazon tanítottuk, hogy a negatív osztály teljes változatosságát
kihasználjuk. Az adathalmazt rétegzett (stratifikált) módon, rögzített
véletlenmaggal ($seed = 42$) osztottuk fel tanító, validációs és teszthalmazra
$70{:}15{:}15$ arányban; a teszthalmazt a tanítás teljes folyamata alatt
érintetlenül hagytuk, és kizárólag a végső kiértékeléshez használtuk fel. A
felosztás eredményét a \ref{tab:split}.~táblázat foglalja össze.

\begin{table}[H]
  \centering
  \caption{Az ISIC 2024 adathalmaz rétegzett felosztása (teljes adathalmaz, $seed=42$).}
  \label{tab:split}
  \begin{tabular}{lrr}
    \hline
    \textbf{Részhalmaz} & \textbf{Összes kép} & \textbf{Rosszindulatú (pozitív)} \\
    \hline
    Tanító (train)      & 280\,741 & 275 \\
    Validációs (val)    & 60\,159  & 59  \\
    Teszt (test)        & 60\,159  & 59  \\
    \hline
    \textbf{Összesen}   & 401\,059 & 393 \\
    \hline
  \end{tabular}
\end{table}

% =============================================================================
\section{A Vision Transformer modell}
% =============================================================================

A kutatás modellfejlesztési megközelítése a \textbf{transzfer tanulás} (transfer
learning): az ImageNet-21k adathalmazon előtanított, általános vizuális
reprezentációkat hordozó súlyokat a bőrelváltozás-osztályozás szűk doménjére
hangoltuk finomhangolással (fine-tuning). Ez a stratégia különösen indokolt, ha
a céldoméni adathalmaz korlátozott számú pozitív példát tartalmaz -- az ISIC 2024
esetén mindössze 393 rosszindulatú esetet a közel 401~000 képből --, mivel az
előtanított jellemzőkivonat (textúrák, alakok, szélek) hatékonyan transzferálható
még az orvosi képi feladatokra is \cite{raghu2019transfusion}.

A fő modell a \texttt{vit\_base\_patch16\_224} architektúra
\cite{dosovitskiy2021vit}, amelyet a \texttt{timm} könyvtárból
\cite{wightman2019timm} töltöttünk be. A törzsháló a
$224 \times 224$ pixeles bemenetből egy 768 dimenziós reprezentációt állít elő,
amelyre egy egyedi osztályozó fejet illesztettünk: rétegnormalizálás, dropout
($p=0{,}3$), egy 256 neuronos lineáris réteg GELU aktivációval, majd egy újabb
dropout ($p=0{,}2$) és a kétosztályos kimeneti réteg.

\subsection{Kétfázisú tanítás}

A modellt két fázisban tanítottuk, hogy egyszerre kezeljük az
osztály-egyensúlytalanságot és a kimeneti valószínűségek kalibrációját:

\begin{itemize}
  \item \textbf{1. fázis -- kiegyensúlyozott tanítás.} Súlyozott
  mintavételezéssel (\texttt{WeightedRandomSampler}) közel $1{:}1$ arányú
  batch-eket állítottunk elő, így a ritka pozitív esetek a tanítás során
  felülreprezentálttá váltak. A veszteségfüggvény súlyozott fókuszveszteség
  \cite{lin2017focal} volt.
  \item \textbf{2. fázis -- kalibráció.} A modellt a valós, erősen
  kiegyensúlyozatlan eloszlású adaton, tizedére csökkentett tanulási rátával
  tanítottuk tovább, hogy a kimeneti valószínűségek a tényleges priorhoz
  igazodjanak.
\end{itemize}

\subsection{Az osztály-egyensúlytalanság kezelése}

Az egyensúlytalanságot három, egymást kiegészítő technikával kezeltük: (i)
súlyozott mintavételezés a kiegyensúlyozott batch-ekhez; (ii) súlyozott
fókuszveszteség, amely a nehéz, félreosztályozott eseteket helyezi előtérbe
(\ref{eq:focal}. egyenlet, $\alpha = 0{,}25$, $\gamma = 2{,}0$, a pozitív
osztály súlya $7{,}0$); valamint (iii) a kisebbségi osztályra alkalmazott
erősebb képaugmentáció. A számítási költség kézben tartására az 1. fázis
epochonkénti mintaszámát 50~000-re korlátoztuk; mivel a kiegyensúlyozott
mintavételezés a 275 pozitív esetet erősen felülmintázza, ez a korlátozás a
tanítási stratégiát érdemben nem befolyásolja, viszont gyakoribb validációt és
korai leállítást tesz lehetővé.

\subsection{Adataugmentáció}

A túltanulás csökkentésére tanítás közben véletlenszerű augmentációt
alkalmaztunk. A többségi (jóindulatú) osztályra alapszintű transzformációkat --
vízszintes és függőleges tükrözés, elforgatás, szín-jitter és affin torzítás --,
a kisebbségi (rosszindulatú) osztályra ezeknél erősebb augmentációt (nagyobb
szögű forgatás, erősebb szín-jitter, perspektivikus torzítás, elmosás)
használtunk, hogy az egyes pozitív példák minél több, eltérő változatban
kerüljenek a modell elé. A validációs és teszthalmazon nem alkalmaztunk
augmentációt, biztosítva a determinisztikus, reprodukálható kiértékelést.

\subsection{Hiperparaméterek}

A tanítás hiperparamétereit a \ref{tab:hyperparams}.~táblázat foglalja össze. Az
optimalizáló AdamW volt, differenciált tanulási rátával (a törzsháló kisebb, a
fej nagyobb rátát kapott), háromepochos lineáris bemelegítéssel, amelyet koszinusz
ütemezésű csökkentés követett. A tanítás vegyes pontosságú (mixed precision)
módban, gradiens-akkumulációval futott, az effektív batch-méret 128 volt.

\begin{table}[H]
  \centering
  \caption{A ViT modell tanításának fő hiperparaméterei.}
  \label{tab:hyperparams}
  \begin{tabular}{ll}
    \hline
    \textbf{Paraméter} & \textbf{Érték} \\
    \hline
    Architektúra            & \texttt{vit\_base\_patch16\_224} (ImageNet-előtanított) \\
    Bemeneti felbontás      & $224 \times 224$ \\
    Batch-méret             & 32 (effektív 128, grad.\ akkumuláció $\times 4$) \\
    Optimalizáló            & AdamW, súlycsökkenés $10^{-2}$ \\
    Tanulási ráta (törzs)   & $2 \times 10^{-6}$ \\
    Tanulási ráta (fej)     & $2 \times 10^{-5}$ \\
    Ütemezés                & 3 epoch warmup + koszinusz csökkenés \\
    1.\ fázis epochszám     & 18 (korai leállítás, türelem $=5$) \\
    2.\ fázis epochszám     & 3 (kalibráció) \\
    Mintaszám / epoch       & 50\,000 \\
    Fókuszveszteség         & $\alpha = 0{,}25$, $\gamma = 2{,}0$, pozitív súly $= 7{,}0$ \\
    \hline
  \end{tabular}
\end{table}

% =============================================================================
\section{A ResNet-18 viszonyítási alap}
% =============================================================================

Viszonyítási alapként egy szintén transzfer tanulással adaptált, ImageNet-en
előtanított ResNet-18 hálót \cite{he2016resnet} finomhangoltunk, amely a mély
tanulásban széles körben alkalmazott, standard konvolúciós referenciamodell.
A ResNet-18 \textbf{pontosan ugyanazt} a feldolgozási folyamatot használja, mint
a ViT -- azonos adatfelosztás
(azonos véletlenmaggal, így azonos teszthalmaz), azonos egyensúlytalanság-kezelés,
azonos kiértékelési protokoll --, így a teljesítménykülönbség értelmezésekor a
tanítási és kiértékelési feltételek nem jelentenek zavaró tényezőt. A
finomhangoláshoz -- mivel a konvolúciós hálók nagyobb tanulási rátát viselnek el
-- a törzsháló, illetve a fej tanulási rátáját $1 \times 10^{-5}$, illetve
$1 \times 10^{-4}$ értékre emeltük; minden egyéb beállítás megegyezett a ViT-ével.

Fontos ugyanakkor megjegyezni, hogy a ResNet-18 lényegesen kevesebb paraméterrel
rendelkezik (mintegy 11~millió), mint a ViT-Base (mintegy 86~millió). A két
modell tehát nem csupán az architektúrában (konvolúció kontra figyelem), hanem a
kapacitásban is különbözik; a ResNet-18-at ezért \emph{standard viszonyítási
alapként} kezeljük, nem pedig szigorúan azonos kapacitású kontrollként. A
megfigyelt teljesítménykülönbséget ennek megfelelően a két tényező együttes
hatásaként értelmezzük.

% =============================================================================
\section{A Gemini zero-shot viszonyítási alap}
% =============================================================================

Második viszonyítási alapként a Google Gemini multimodális nagy nyelvi modelljét
\cite{gemini2023} (a \texttt{gemini-2.5-flash} változatot) értékeltük
\textit{zero-shot} módban, azaz célzott betanítás nélkül. A modell minden képhez
egy minimális szöveges utasítást kapott, amely a rosszindulatúság
valószínűségének $[0,1]$ skálán történő becslésére kérte; sem példákat
(few-shot), sem klinikai metaadatot nem adtunk a bemenethez. Mivel a Gemini
API-alapú, soronkénti lekérdezése a teljes, mintegy 60~000 elemű teszthalmazon
költséges és lassú lenne, a Gemini-kiértékelést egy 2~000 képes részhalmazon
végeztük, amely a teszthalmaz \textbf{összes} pozitív esetét és véletlenszerűen
mintavételezett negatív eseteket tartalmazott; a tisztességes összehasonlítás
érdekében a ViT modellt ugyanezen a részhalmazon is kiértékeltük.

A Gemini szerepe a dolgozatban kettős szempontból tanulságos: egyrészt
megmutatja, mire képes egy általános célú, célzott betanítás nélküli multimodális
modell ezen a feladaton, másrészt rávilágít a célzottan betanított, kalibrált
osztályozó modell előnyeire -- összhangban a szakirodalmi megfigyelésekkel
\cite{cirone2024llm}.

% =============================================================================
\section{Kísérleti elrendezés}
% =============================================================================

A modelleket a Kaggle felhőplatformján, NVIDIA Tesla T4 GPU-n tanítottuk és
értékeltük. Az implementáció a PyTorch keretrendszerre \cite{paszke2019pytorch}, a
modellek betöltése a \texttt{timm} könyvtárra \cite{wightman2019timm}, a metrikák
és statisztikai számítások pedig a scikit-learn könyvtárra
\cite{pedregosa2011sklearn} épültek. A reprodukálhatóság érdekében minden
véletlenszerű folyamatot rögzített maggal ($seed = 42$) inicializáltunk.

A kiértékelést két klinikailag értelmezhető \textit{működési ponton} végeztük: a
$\ge 0{,}90$ érzékenységnél (klinikai szűrési standard) és a $\ge 0{,}80$
érzékenységnél (az ISIC 2024 verseny pAUC-küszöbe). A döntési küszöböt minden
esetben a validációs halmazon kalibráltuk, majd a rögzített küszöböt alkalmaztuk
a teszthalmazon. A teljesítményt a \ref{eq:sens-spec}--\ref{eq:ppv-npv}.
egyenletekben definiált metrikákkal, az AUC-val és a hivatalos ISIC 2024 pAUC-val
mértük; a pontbecslések bizonytalanságát Clopper--Pearson-féle (arányok) és
DeLong-féle (AUC) konfidencia-intervallumokkal \cite{sun2014delong} jelöltük, a
modellek közötti AUC-különbségeket pedig a DeLong-féle páros teszttel
vizsgáltuk. A valószínűségi kalibrációt a Brier-pontszámmal
\cite{guo2017calibration} jellemeztük.

% =============================================================================
\section{Eredmények}
% =============================================================================

Ebben a szakaszban a ViT modell teljesítményét értékeljük a teszthalmazon, majd
összevetjük a ResNet-18, az EVA02 és a Gemini viszonyítási alapokkal.

\subsection{A tanítás konvergenciája}

A \ref{fig:training-curves}.~ábra a tanítás során mért veszteség, valamint a
validációs AUC és pAUC alakulását mutatja epochonként. A validációs AUC a kezdeti
$0{,}73$-ról $0{,}90$-re, a pAUC $0{,}055$-ről $0{,}132$-re emelkedett; a legjobb
validációs pAUC a 11.~epochnál állt elő, a korai leállítás pedig a 16.~epochnál
aktiválódott. A második (kalibrációs) fázis a validációs metrikákat érdemben nem
változtatta, így a végső modell lényegében az első fázis legjobb állapota.

\begin{figure}[H]
  \centering
  \includegraphics[width=\textwidth]{images/results/vit_training_curves.png}
  \caption{A ViT modell tanításának konvergenciája: a tanítási veszteség,
    valamint a validációs AUC és pAUC alakulása az epochok függvényében.}
  \label{fig:training-curves}
\end{figure}

\subsection{A ViT teljesítménye a teszthalmazon}

A \ref{tab:vit-results}.~táblázat a ViT modell teljesítményét foglalja össze a
két működési ponton.

\begin{table}[H]
  \centering
  \caption{A ViT modell teljesítménye a teszthalmazon (60\,159 kép, 59 pozitív), két működési ponton.}
  \label{tab:vit-results}
  \begin{tabular}{lcc}
    \hline
    \textbf{Metrika} & \textbf{Érzékenység $\ge 0{,}90$} & \textbf{Érzékenység $\ge 0{,}80$} \\
    \hline
    AUC (95\% CI)          & \multicolumn{2}{c}{$0{,}939$ ($0{,}915$--$0{,}963$)} \\
    pAUC (TPR $\ge 0{,}80$) & \multicolumn{2}{c}{$0{,}159$} \\
    Brier-pontszám         & \multicolumn{2}{c}{$0{,}047$} \\
    \hline
    Küszöb                 & $0{,}170$ & $0{,}270$ \\
    Érzékenység (95\% CI)  & $0{,}966$ ($0{,}883$--$0{,}996$) & $0{,}881$ ($0{,}771$--$0{,}951$) \\
    Specificitás (95\% CI) & $0{,}668$ ($0{,}664$--$0{,}671$) & $0{,}843$ ($0{,}840$--$0{,}846$) \\
    PPV / NPV              & $0{,}003$ / $1{,}000$ & $0{,}005$ / $1{,}000$ \\
    \hline
  \end{tabular}
\end{table}

A \ref{fig:roc}.~ábra a ViT modell ROC-görbéjét mutatja, kiemelve a pAUC
számításához használt, $TPR \ge 0{,}80$ tartományt.

\begin{figure}[H]
  \centering
  \includegraphics[width=0.7\textwidth]{images/results/vit_roc_curve.png}
  \caption{A ViT modell ROC-görbéje a teszthalmazon, a pAUC számításához használt
    $TPR \ge 0{,}80$ tartomány kiemelésével.}
  \label{fig:roc}
\end{figure}

\subsection{Összehasonlítás a viszonyítási alapokkal}

A \ref{tab:comparison}.~táblázat a négy modellt veti össze a két közvetlenül
összehasonlítható minőség-mérőszám -- az AUC és a hivatalos ISIC 2024 pAUC --
mentén. Mindkét mutató a modellek rangsorolási képességét méri, és független a
döntési küszöbtől, így akkor is értelmes az összevetésük, ha az egyes modellek
saját, szivárgásmentes (held-out) halmazukon kerültek kiértékelésre.

\begin{table}[H]
  \centering
  \caption{A négy modell összehasonlítása az AUC és a hivatalos pAUC mentén. Az
    EVA02 a saját, szivárgásmentes validációs foldján szerepel (a modellt a teljes
    adat 80\%-án tanították, így a mi teszthalmazunkon történő kiértékelése
    adatszivárgást okozna); a Gemini csak egy 2\,000 elemű részhalmazon futott, a
    pAUC-ja a korábbi (eltérő) definícióval készült, ezért nem szerepel.}
  \label{tab:comparison}
  \begin{tabular}{lccl}
    \hline
    \textbf{Modell} & \textbf{AUC} & \textbf{pAUC} & \textbf{Kiértékelési halmaz} \\
    \hline
    ViT-Base (saját, finomhangolt)     & $0{,}939$ & $0{,}159$ & teszt (60\,159, 59 poz.) \\
    EVA02 (legjobb image-only ref.)    & $0{,}929$ & $0{,}158$ & val-fold (1\,651, 79 poz.) \\
    ResNet-18 (CNN, viszonyítási alap) & $0{,}876$ & $0{,}122$ & teszt (60\,159, 59 poz.) \\
    Gemini 2.5 Flash (zero-shot)       & $0{,}740$ & --        & részhalmaz (2\,000, 59 poz.) \\
    \hline
  \end{tabular}
\end{table}

A \ref{fig:model-comparison}.~ábra a négy modell AUC- és pAUC-értékeit
szemlélteti.

\begin{figure}[H]
  \centering
  \includegraphics[width=\textwidth]{images/diagrams/model_comparison.png}
  \caption{A négy modell összehasonlítása az AUC (bal) és a hivatalos ISIC 2024
    pAUC (jobb) mentén. A Gemini pAUC-értéke nem szerepel, mivel az eltérő
    definícióval, kisebb részhalmazon készült. A pAUC elméleti maximuma $0{,}20$.}
  \label{fig:model-comparison}
\end{figure}

\paragraph{ViT kontra ResNet-18 (azonos teszthalmaz).} A két modell pontosan
ugyanazon a teszthalmazon, azonos feldolgozási folyamattal került kiértékelésre,
így a tanítási és kiértékelési feltételek nem zavaró tényezők (a kapacitásbeli
különbségre az előző szakaszban tett megjegyzés mellett). A ViT minden mutatóban
érdemben felülmúlja a ResNet-et: AUC $0{,}939$ kontra $0{,}876$,
pAUC $0{,}159$ kontra $0{,}122$. A
különbség a magas érzékenységű tartományban a legmarkánsabb: a $\ge 0{,}80$
érzékenységi ponton a specificitás $0{,}843$ kontra $0{,}625$, a $\ge 0{,}90$
ponton pedig $0{,}668$ kontra $0{,}400$ -- vagyis a ViT azonos érzékenység
mellett lényegesen kevesebb téves riasztást ad.

\paragraph{ViT kontra EVA02 (a legjobb image-only Kaggle modell).} A ViT-Base
gyakorlatilag egyenrangú a verseny legerősebb tisztán képi modelljével: a
hivatalos pAUC szinte azonos ($0{,}159$ kontra $0{,}158$), az AUC pedig
hasonló ($0{,}939$ kontra $0{,}929$). Mivel mindkettő szivárgásmentes halmazon, a
hivatalos metrikával készült, a párhuzam tisztességes -- noha nem ugyanazon a
halmazon. Fontos azonban, hogy az EVA02 a natív $336\times336$ felbontásán fut
(szemben a ViT $224\times224$-ével), ezért ez nem szigorú, azonos felbontású
kontroll, hanem a legjobb publikált image-only megoldással szembeni viszonyítás.
A két modell specificitása nem vethető össze közvetlenül, mivel az EVA02
kiértékelési halmaza jóval kisebb és eltérő prevalenciájú.

\paragraph{ViT kontra Gemini (zero-shot).} A célzott betanítás nélküli, általános
célú Gemini modell jelentősen elmarad (AUC $0{,}740$): hajlamos a legtöbb
elváltozást gyanúsnak jelölni, ezért a specificitása alacsony. Ez megerősíti a
rendszer tervezési döntését, amely a klasszifikációt egy célzottan betanított
modellre, a magyarázatot pedig a nagy nyelvi modellre bízza.

\subsection{Statisztikai elemzés}

A modellek rangsorolási teljesítményének bizonytalanságát a DeLong-féle 95\%-os
konfidencia-intervallumok jellemzik. A ViT AUC-jának intervalluma
($0{,}915$--$0{,}963$) és a ResNet-é ($0{,}837$--$0{,}914$) gyakorlatilag nem
fed át, ami a két modell közötti különbség statisztikai szignifikanciájára utal.
A mentett predikciókból a DeLong-féle páros teszt formálisan is elvégezhető, mivel
a ViT és a ResNet ugyanazon a teszthalmazon került kiértékelésre.

A valószínűségi kalibrációt a Brier-pontszám jellemzi. Mindkét képi modell
Brier-pontszáma (ViT: $0{,}047$, ResNet: $0{,}075$) magasabb a triviális, a
prevalenciát jósló alapvonalénál ($0{,}001$), ami a kimeneti valószínűségek
gyenge kalibráltságát jelzi -- ez az erős osztály-egyensúlytalanság és a
kiegyensúlyozott mintavételezés ismert következménye. A ViT ugyanakkor jobban
kalibrált, mint a ResNet. Hangsúlyozandó, hogy a \textit{rangsorolás} (AUC, pAUC)
erős; a nyers valószínűségek utólagos kalibrálása a jövőbeli munka tárgya.

\subsection{Az eredmények megbeszélése és korlátai}

Az eredmények fő tanulsága, hogy a finomhangolt ViT-Base modell eléri a verseny
legjobb tisztán képi megoldásának (EVA02) szintjét, miközben jelentősen felülmúlja
mind a klasszikus konvolúciós alapot (ResNet-18), mind a zero-shot nagy nyelvi
modellt (Gemini). A magas érzékenységű működési ponton mért közel tökéletes
negatív prediktív érték (NPV $\approx 1{,}0$) klinikai szűrőeszközként kedvező: a
,,jó\-indulatú'' besorolás megbízható.

Az értelmezésnél több korlátot figyelembe kell venni. \emph{(i)} A teszthalmazban
csak 59 pozitív eset szerepel, ami a metrikák -- különösen az érzékenység és a
pAUC -- becslését érzékennyé teszi a mintavételi zajra; ezt a megadott
konfidencia-intervallumok tükrözik. \emph{(ii)} A pozitív prediktív érték (PPV)
nagyon alacsony ($\approx 0{,}005$), ami az extrém, $\sim 1{:}1000$ arányú
osztály-egyensúlytalanság közvetlen következménye -- nagyszámú negatív eset
mellett még jó specificitás mellett is sok a téves pozitív. \emph{(iii)} Az EVA02
eltérő (kisebb, magasabb prevalenciájú) halmazon és felbontáson került
kiértékelésre, ezért referenciaként, nem szigorú kontrollként értelmezendő.
\emph{(iv)} A kimeneti valószínűségek kalibrációja további javítást igényel.
